\documentclass[a4paper, 12pt]{article}

\usepackage[portuges]{babel}
\usepackage[utf8]{inputenc}
\usepackage{amsmath}
\usepackage{indentfirst}
\usepackage{graphicx}
\usepackage{multicol,lipsum}
\usepackage{geometry}
\usepackage{amsmath, amssymb}
\usepackage{graphicx}
\usepackage{hyperref}
\usepackage{listings}
\usepackage{xcolor}
\usepackage{booktabs}
\usepackage{longtable}
\usepackage{caption}
\usepackage{float}

\begin{document}
%\maketitle

\begin{titlepage}
	\begin{center}
	
	%\begin{figure}[!ht]
	%\centering
	%\includegraphics[width=2cm]{c:/ufba.jpg}
	%\end{figure}

		\Huge{Universidade Federal do Cariri}\\
		\large{Ciência da Computação}\\ 
		\large{Lógica}\\ 
		\vspace{15pt}
        \vspace{95pt}
        \textbf{\LARGE{Chaos Sudoku: Formalização em Lógica Proposicional e Verificação Experimental com SAT Solver}}\\
		%\title{{\large{Título}}}
		\vspace{3,5cm}
	\end{center}
	
	\begin{flushleft}
		\begin{tabbing}
			Alunos: David Hudson Gomes Alves e Vitória do Nascimento Pontes \\
			Professor: Luis Henrique Bustamante de Morais \\
	\end{tabbing}
 \end{flushleft}
	\vspace{1cm}
	
	\begin{center}
		\vspace{\fill}
			 Agosto\\
		 2026
			\end{center}
\end{titlepage}
%%%%%%%%%%%%%%%%%%%%%%%%%%%%%%%%%%%%%%%%%%%%%%%%%%%%%%%%%%%


% % % % % % % % % % % % % % % % % % % % % % % % % %
\newpage
\tableofcontents
\thispagestyle{empty}

\newpage
\pagenumbering{arabic}
% % % % % % % % % % % % % % % % % % % % % % % % % % %

\section{Introdução}

O trabalho proposto tem como finalidade aproximar a teoria da Lógica Proposicional de uma aplicação computacional concreta. O enunciado solicita a escolha de um problema combinatório, sua formalização em lógica proposicional, a geração de instâncias no formato DIMACS e a utilização de um SAT solver moderno para verificar experimentalmente a satisfatibilidade.

Neste projeto foi escolhido o \emph{Chaos Sudoku}, uma variação do Sudoku na qual as regiões podem ser irregulares. Cada possível colocação de um valor em uma célula pode ser tratada como uma proposição, e as regras do jogo podem ser expressas por restrições lógicas. A utilização de SAT é particularmente interessante porque o solver não precisa conhecer as regras específicas do Sudoku. Ele recebe somente uma fórmula CNF. Assim, o conhecimento do domínio fica concentrado na etapa de modelagem, enquanto o processo de busca por uma atribuição satisfatória é delegado ao solver.

\section{Descrição do problema e das instâncias}

Uma instância de Chaos Sudoku é composta por uma grade quadrada \(N \times N\). Cada célula deve receber um inteiro entre 1 e \(N\). Além da restrição de preenchimento, cada valor deve aparecer exatamente uma vez em cada linha, exatamente uma vez em cada coluna e exatamente uma vez em cada região irregular. As regiões são fornecidas explicitamente por uma matriz \(N \times N\). O identificador presente em cada posição informa a região à qual aquela célula pertence. Para uma instância válida, existem exatamente \(N\) regiões, identificadas por \(1,\dots,N\), e cada região possui exatamente \(N\) células.

\subsection{Representação JSON}

O projeto representa cada instância por duas matrizes: regiões e pistas. Um exemplo \(4\times4\) é:

\begin{verbatim}
{
  "regioes": [
    [3, 1, 1, 1],
    [3, 3, 3, 1],
    [2, 4, 4, 4],
    [2, 2, 2, 4]
  ],
  "pistas": [
    [1, 0, 0, 0],
    [0, 3, 0, 0],
    [0, 0, 1, 0],
    [0, 0, 0, 3]
  ]
}
\end{verbatim}

Na matriz de pistas, o valor zero significa que a célula está inicialmente vazia. Um valor entre 1 e \(N\) fixa aquele valor como parte da instância.

\subsection{Diferença em relação ao Sudoku tradicional}

No Sudoku tradicional, as regiões normalmente possuem uma forma regular e sua geometria pode ser inferida a partir de \(N\). No Chaos Sudoku, essa hipótese não é feita. A matriz regiões é parte da entrada e permite descrever regiões irregulares. Isso torna a formalização mais geral: a mesma lógica de unicidade é aplicada ao conjunto de células que compartilham um identificador de região.

\subsection{Instâncias utilizadas}

O repositório contém sete fórmulas CNF correspondentes às instâncias \texttt{exemplo\_4}, \texttt{exemplo\_4\_unsat}, \texttt{exemplo\_5}, \texttt{exemplo\_6}, \texttt{exemplo\_7}, \texttt{exemplo\_8} e \texttt{exemplo\_9}. Os tamanhos crescem de \(4\times4\) até \(9\times9\), permitindo observar o aumento da representação.

\section{Formalização matemática}

\subsection{Variáveis proposicionais}

Seja \(N\) o tamanho da grade. Definimos o conjunto de variáveis:
\[
X = \{ x_{ijv} \mid 1 \leq i, j, v \leq N \}.
\]
A interpretação é direta: \(x_{ijv}\) é verdadeira se, e somente se, a célula localizada na linha \(i\) e coluna \(j\) contém o valor \(v\). Para uma célula existem \(N\) possibilidades e, portanto, são criadas \(N\) variáveis.

Como existem \(N^2\) células, o número total de variáveis é:
\[
|X| = N^3.
\]
Essa escolha é conhecida como codificação \emph{one-hot}: em vez de representar diretamente o número escrito na célula, representamos separadamente cada possibilidade de valor.

\subsection{Codificação dos identificadores DIMACS}

O DIMACS utiliza inteiros positivos para representar variáveis. O projeto associa cada \(x_{ijv}\) ao identificador:
\[
\text{var}(i,j,v) = N^2(i-1) + N(j-1) + v.
\]
A primeira parcela pula todas as linhas anteriores; a segunda pula as células anteriores da linha atual; a terceira seleciona o valor dentro da célula. O maior identificador possível é \(N^3\), exatamente a quantidade de variáveis.

Por exemplo, para \(N=4\), a célula \((1,1)\) usa os IDs \(1,2,3,4\); a célula \((1,2)\) usa \(5,6,7,8\); e a célula \((2,1)\) começa no ID 17. Isso permite converter a estrutura tridimensional dos índices \((i,j,v)\) para a numeração linear exigida pelo DIMACS.

\subsection{Família F1: pelo menos um valor}

Para cada célula \((i,j)\), é necessário garantir que algum valor seja escolhido:
\[
(x_{ij1} \lor x_{ij2} \lor \dots \lor x_{ijN}).
\]
Existe uma cláusula desse tipo para cada uma das \(N^2\) células. Logo, F1 possui \(N^2\) cláusulas. Essa família, isoladamente, permite que mais de um valor seja escolhido na mesma célula; por isso é necessária a família seguinte.

\subsection{Família F2: no máximo um valor por célula}

Para impedir que uma célula contenha dois valores simultaneamente, para cada par de valores \(v_1 < v_2\) incluímos:
\[
(\neg x_{ijv_1} \lor \neg x_{ijv_2}).
\]
Para cada célula existem \(\binom{N}{2}\) pares. Como existem \(N^2\) células, a quantidade de cláusulas de F2 é \(N^2 \cdot \binom{N}{2}\). Com F1 e F2 em conjunto, cada célula recebe exatamente um valor.

\subsection{Família F3: unicidade nas linhas}

Para cada linha \(i\) e valor \(v\), o mesmo valor não pode aparecer em duas colunas distintas. Para cada par \(j_1 < j_2\):
\[
(\neg x_{ij_1v} \lor \neg x_{ij_2v}).
\]
Há \(N\) linhas, \(N\) valores e \(\binom{N}{2}\) pares de colunas, totalizando \(N^2 \cdot \binom{N}{2}\) cláusulas.

\subsection{Família F4: unicidade nas colunas}

Analogamente, para cada coluna \(j\) e valor \(v\), duas linhas diferentes não podem receber o mesmo valor. Para \(i_1 < i_2\):
\[
(\neg x_{i_1jv} \lor \neg x_{i_2jv}).
\]
Também são geradas \(N^2 \cdot \binom{N}{2}\) cláusulas.

\subsection{Família F5: unicidade nas regiões}

Se duas células pertencem à mesma região \(r\), elas não podem conter o mesmo valor. Para cada região, valor e par de células distintas da região:
\[
(\neg x_{i_1j_1v} \lor \neg x_{i_2j_2v}).
\]
Como cada instância válida possui \(N\) regiões e cada região contém exatamente \(N\) células, cada região também gera \(\binom{N}{2}\) pares para cada valor. Portanto, F5 possui \(N^2 \cdot \binom{N}{2}\) cláusulas.

\subsection{Família F6: pistas}

Uma pista \((i,j)=v\) é representada por uma cláusula unitária:
\[
x_{ijv}.
\]
Se a instância possui \(G\) pistas, F6 acrescenta exatamente \(G\) cláusulas. Uma pista contraditória não precisa ser detectada pelo gerador: ela pode ser representada normalmente e fará com que a fórmula seja UNSAT.

\subsection{Contagem total}

Somando as famílias:
\[
M = N^2 + 4N^2 \cdot \binom{N}{2} + G.
\]
Como \(\binom{N}{2} = \frac{N(N-1)}{2}\), obtemos:
\[
M = N^2 \bigl(1 + 2N(N-1)\bigr) + G.
\]
Assim, a fórmula possui \(N^3\) variáveis e \(M\) cláusulas. O termo dominante cresce proporcionalmente a \(N^4\).

\section{Implementação do gerador DIMACS}

O gerador está implementado em \texttt{src/gerador.py} e utiliza somente recursos da biblioteca padrão do Python. A interface recebe dois argumentos principais: o tamanho \(N\) e o caminho do arquivo JSON.

\begin{verbatim}
python3 src/gerador.py N ARQUIVO_JSON
\end{verbatim}

A saída é produzida na saída padrão. Isso permite redirecionar diretamente o resultado para um arquivo \texttt{.cnf}.

\begin{verbatim}
python3 src/gerador.py 4 instances/exemplo_4.json > exemplo_4.cnf
\end{verbatim}

\subsection{Leitura e validação da instância}

A função \texttt{ler\_instancia} abre o JSON e delega a validação estrutural para \texttt{validar\_instancia}. O gerador verifica a existência do campo \texttt{regioes}, o formato \(N\times N\) das matrizes, os identificadores de regiões e a quantidade de células por região. O campo \texttt{pistas} é opcional; quando não existe, todas as células são tratadas como vazias.

É importante observar que essas verificações são estruturais. O programa não tenta decidir previamente se as pistas são compatíveis com uma solução. Essa decisão é justamente uma das funções do SAT solver.

\subsection{Organização interna das cláusulas}

Cada família de restrição é construída por uma função específica. As cláusulas são representadas internamente por listas de inteiros, em que um inteiro positivo representa uma variável e um inteiro negativo representa sua negação.

\begin{verbatim}
def var(i, j, v, n):
    return n**2 * (i - 1) + n * (j - 1) + v
\end{verbatim}

Para gerar pares sem repetição, o código utiliza \texttt{combinations(..., 2)}, evitando a criação duplicada de uma mesma cláusula. A estrutura de regiões é obtida a partir da matriz fornecida pela instância.

\subsection{Geração do DIMACS}

O cabeçalho DIMACS segue o formato \texttt{'p cnf N M'}. No primeiro exemplo, o arquivo começa com \texttt{'p cnf 64 404'}: 64 representa \(N^3\) para \(N=4\) e 404 corresponde às 400 cláusulas estruturais mais quatro pistas.

Cada cláusula é encerrada pelo literal 0. Por exemplo, a cláusula de pelo menos um valor para a primeira célula de uma instância \(4\times4\) é escrita como:
\begin{verbatim}
1 2 3 4 0
\end{verbatim}

\section{Experimentação com o SAT Solver}

A etapa experimental utiliza o CaDiCaL, conforme recomendado no enunciado da disciplina. O objetivo é fornecer ao solver as fórmulas CNF geradas e registrar o resultado de satisfatibilidade, o tamanho da fórmula e o tempo de execução. Foram utilizadas sete instâncias, cobrindo tamanhos de \(4\times4\), \(5\times5\), \(6\times6\), \(7\times7\), \(8\times8\) e \(9\times9\). Os arquivos CNF presentes no projeto registram os seguintes cabeçalhos:

\begin{table}[H]
\centering
\caption{Resultados experimentais com CaDiCaL.}
\begin{tabular}{c r r r l r r}
\toprule
Instância & N & Variáveis & Cláusulas & Pistas & Resultado & Tempo (s) \\
\midrule
exemplo\_4 & 4 & 64 & 404 & 4 & SAT & 0,013 \\
exemplo\_4\_unsat & 4 & 64 & 402 & 2 & UNSAT & 0,001 \\
exemplo\_5 & 5 & 125 & 1027 & 2 & UNSAT & 0,002 \\
exemplo\_6 & 6 & 216 & 2202 & 6 & SAT & 0,003 \\
exemplo\_7 & 7 & 343 & 4167 & 2 & UNSAT & 0,002 \\
exemplo\_8 & 8 & 512 & 7240 & 8 & SAT & 0,003 \\
exemplo\_9 & 9 & 729 & 11754 & 9 & SAT & 0,003 \\
\bottomrule
\end{tabular}
\end{table}

\subsection{Instância exemplo\_4}

A instância \(4\times4\) possui 64 variáveis, 404 cláusulas e quatro pistas. O resultado é SAT. Uma solução reconstruída a partir do modelo é:
\[
\begin{array}{cccc}
1 & 4 & 3 & 2 \\
4 & 3 & 2 & 1 \\
3 & 2 & 1 & 4 \\
2 & 1 & 4 & 3
\end{array}
\]
Essa grade satisfaz as quatro condições do problema: cada linha contém \(1,2,3,4\) sem repetição; o mesmo ocorre nas colunas; e cada região contém os quatro valores.

\subsection{Instância exemplo\_4\_unsat}

Essa instância mantém as mesmas regiões, mas possui duas pistas iguais a 1 na primeira linha. As pistas obrigam simultaneamente \(x_{111}\) e \(x_{121}\) a serem verdadeiras. Entretanto, F3 contém a cláusula \((\neg x_{111} \lor \neg x_{121})\), produzindo uma contradição direta. O solver, portanto, retorna UNSATISFIABLE.

\subsection{Instância exemplo\_5}

A instância \(5\times5\) possui 125 variáveis e 1.027 cláusulas. Ela contém duas pistas iguais a 1 na primeira linha. Pela mesma razão da instância \(4\times4\) insatisfatível, as pistas conflitam com a restrição de unicidade por linha. O resultado registrado é UNSAT.

\subsection{Instância exemplo\_6}

A instância \(6\times6\) possui 216 variáveis, 2.202 cláusulas e seis pistas. As regiões são seis blocos regulares de \(3\times2\) células. O resultado é SAT. A solução reconstruída é:
\[
\begin{array}{cccccc}
1 & 6 & 5 & 3 & 4 & 2 \\
3 & 2 & 4 & 6 & 1 & 5 \\
6 & 4 & 3 & 5 & 2 & 1 \\
5 & 1 & 2 & 4 & 6 & 3 \\
2 & 3 & 6 & 1 & 5 & 4 \\
4 & 5 & 1 & 2 & 3 & 6
\end{array}
\]

\subsection{Instância exemplo\_7}

A instância \(7\times7\) possui 343 variáveis e 4.167 cláusulas. Assim como \texttt{exemplo\_4\_unsat} e \texttt{exemplo\_5}, suas duas primeiras pistas são 1 na mesma linha, tornando a fórmula insatisfatível. O solver retorna UNSAT.

\subsection{Instância exemplo\_8}

A instância \(8\times8\) possui 512 variáveis e 7.240 cláusulas, com oito pistas. As regiões formam oito blocos de \(4\times2\) células. O resultado é SAT e uma solução reconstruída é:
\[
\begin{array}{cccccccc}
1 & 3 & 7 & 5 & 6 & 8 & 4 & 2 \\
8 & 2 & 4 & 6 & 7 & 3 & 5 & 1 \\
7 & 6 & 3 & 2 & 8 & 4 & 1 & 5 \\
5 & 8 & 1 & 4 & 2 & 7 & 3 & 6 \\
6 & 7 & 8 & 3 & 5 & 1 & 2 & 4 \\
4 & 5 & 1 & 3 & 6 & 8 & 7 & 2 \\
4 & 5 & 6 & 8 & 1 & 2 & 7 & 3 \\
3 & 1 & 2 & 7 & 4 & 5 & 6 & 8
\end{array}
\]

\subsection{Instância exemplo\_9}

A maior instância utilizada possui \(N=9\), 729 variáveis e 11.754 cláusulas. As regiões são nove blocos de \(3\times3\) células e existem nove pistas na primeira linha, uma para cada valor de 1 a 9. O resultado é SAT. A grade reconstruída é:
\[
\begin{array}{ccccccccc}
1 & 2 & 3 & 4 & 5 & 6 & 7 & 8 & 9 \\
8 & 7 & 5 & 3 & 9 & 1 & 2 & 4 & 6 \\
6 & 9 & 4 & 8 & 2 & 7 & 1 & 3 & 5 \\
7 & 6 & 8 & 2 & 1 & 5 & 3 & 9 & 4 \\
5 & 3 & 1 & 6 & 4 & 9 & 8 & 2 & 7 \\
2 & 4 & 9 & 7 & 3 & 8 & 5 & 6 & 1 \\
4 & 5 & 6 & 1 & 8 & 2 & 9 & 7 & 3 \\
9 & 8 & 7 & 5 & 6 & 3 & 4 & 1 & 2 \\
3 & 1 & 2 & 9 & 7 & 4 & 6 & 5 & 8
\end{array}
\]

\subsection{Reconstrução do modelo SAT}

A reconstrução é feita interpretando os literais positivos presentes no modelo. Se o solver informa que a variável correspondente a \(x_{ijv}\) é verdadeira, o algoritmo coloca \(v\) na posição \((i,j)\). As famílias F1 e F2 garantem que cada célula tenha exatamente um valor, evitando ambiguidades durante a reconstrução.

Essa etapa é importante porque demonstra que o solver não apenas classifica a fórmula. Em uma instância SAT, seu modelo contém informação suficiente para recuperar uma solução concreta do problema original.

\subsection{Observação sobre os tempos}

Os tempos medidos ficaram entre 1 e 13 milissegundos, sem relação direta com o tamanho da instância: a maior demora observada ocorreu na instância \(4\times4\) (0,013 s), enquanto instâncias maiores, como a \(9\times9\), foram resolvidas em 3 milissegundos. Esse comportamento é esperado para fórmulas desse porte, pequenas o suficiente para que fatores como o tempo de inicialização do solver dominem sobre a dificuldade real de busca. Diferenças de desempenho mais claras entre as instâncias provavelmente só apareceriam em grades bem maiores que \(N=9\).

\section{Análise dos resultados}

\subsection{Crescimento do número de variáveis}

A quantidade de variáveis é \(N^3\). Isso significa que aumentar o tamanho da grade provoca crescimento relativamente rápido da representação. Por exemplo, passar de \(N=4\) para \(N=9\) aumenta o número de variáveis de 64 para 729.

\subsection{Crescimento do número de cláusulas}

O número de cláusulas sem pistas é \(N^2(1+2N(N-1))\). O termo dominante é aproximadamente \(2N^4\). Esse crescimento é mais rápido que o das variáveis porque as restrições de unicidade são codificadas comparando pares de possibilidades.

\begin{table}[H]
\centering
\caption{Crescimento do número de variáveis e cláusulas.}
\begin{tabular}{c r r r}
\toprule
\(N\) & \(N^3\) variáveis & Cláusulas sem pistas & Exemplo no projeto \\
\midrule
4 & 64 & 400 & 404 \\
5 & 125 & 1025 & 1027 \\
6 & 216 & 2196 & 2202 \\
7 & 343 & 4165 & 4167 \\
8 & 512 & 7232 & 7240 \\
9 & 729 & 11745 & 11754 \\
\bottomrule
\end{tabular}
\end{table}

\subsection{Interpretação dos casos UNSAT}

Os casos UNSAT presentes no conjunto experimental são particularmente didáticos. A insatisfatibilidade não decorre de uma falha do gerador, mas da própria instância: duas pistas obrigam o mesmo valor a ocupar duas células da mesma linha. A formalização traduz essa regra em cláusulas e o solver detecta que elas não podem ser satisfeitas em conjunto.

\subsection{Adequação da formalização}

A codificação \emph{one-hot} é adequada ao objetivo do trabalho porque torna a interpretação das variáveis muito simples. Cada variável corresponde diretamente a uma afirmação sobre uma célula. Além disso, as restrições de unicidade podem ser expressas com cláusulas binárias, facilitando a implementação e a explicação matemática.

Como desvantagem, a abordagem não é a mais econômica em número de cláusulas. Para cada conjunto de possibilidades são gerados pares, fazendo com que o tamanho da fórmula cresça aproximadamente como \(N^4\). Uma investigação futura poderia comparar outras codificações ou esquemas de cardinalidade.

\section{Conclusão}

O trabalho demonstrou como um problema combinatório pode ser transformado sistematicamente em uma instância de SAT. A partir de uma grade \(N \times N\), foram definidas \(N^3\) variáveis proposicionais, cada uma representando uma possível atribuição de valor a uma célula. As regras do Chaos Sudoku foram então decompostas em famílias de cláusulas CNF, permitindo uma correspondência clara entre a especificação do problema e a fórmula lógica.

A implementação em Python materializa essa formalização em arquivos DIMACS. A representação de regiões por uma matriz torna o gerador capaz de trabalhar com regiões irregulares, característica central do Chaos Sudoku. A interface por linha de comando também permite integrar o gerador diretamente a um SAT solver.

Os experimentos com sete instâncias mostraram resultados SAT e UNSAT. Nos casos satisfativos, foi possível recuperar as grades a partir dos modelos; nos casos insatisfativos, a causa pode ser explicada pelas pistas conflitantes e pelas cláusulas de unicidade. O crescimento de 64 para 729 variáveis e de 404 para 11.754 cláusulas evidencia como a representação proposicional aumenta com \(N\), embora as instâncias avaliadas continuem pequenas o suficiente para serem resolvidas muito rapidamente.

Dessa forma, a experiência evidencia a principal vantagem da abordagem SAT: depois de formalizado, o problema passa a ser tratado por um mecanismo geral de decisão. A qualidade do resultado depende, entretanto, da escolha da modelagem. Uma continuação natural deste projeto seria comparar a codificação adotada com uma representação alternativa, avaliar outros solvers e estudar técnicas para reduzir o número de cláusulas.

\section{Repositório}

Projeto Chaos Sudoku. Repositório: \url{https://github.com/VitoriaPontes/chaos-sudoku}

\end{document}